\chapter{Потенциал действия: нелинейный триггер}
\label{sec:actionpotential}

\emph{Глава в работе. План:} Уравнения Ходжкина---Хаксли как электрическая схема с двумя нелинейными проводимостями. Положительная обратная связь по натриевому току, порог, рефрактерность и закон «всё или ничего» --- как у триггера Шмитта. Проведение импульса по аксону как регенеративный повторитель, миелин как изоляция между повторителями.
